\title[Undershoot of $h$-fold sumsets]{Undershoot of $h$-fold sumsets: Ap\'ery spectra, similar ordering and the three-distance theorem}

\author*[1]{\fnm{Prashant} \sur{Chaudhary}}\email{prashant@prashantchaudhary.com}
\author[2]{\fnm{Kapil} \sur{Khandelwal}}\email{kapil.khandelwal@nd.edu}

\affil*[1]{\orgname{Independent Researcher}, \orgaddress{\city{San Francisco Bay Area}, \state{CA}, \country{USA}}}
\affil[2]{\orgdiv{Department of Civil and Environmental Engineering and Earth Sciences}, \orgname{University of Notre Dame}, \orgaddress{\city{Notre Dame}, \state{IN}, \country{USA}}}

\abstract{For a finite set $A_0\subset\mathbb N$ with $0\in A_0$, $\max A_0=d$ and $\gcd(A_0)=1$, and a target $T$, the undershoot at level $h$ is the distance from $T$ down to the largest element of the $h$-fold sumset $hA_0$ not above $T$. Its total over $T\in[0,hd]$ is bounded in $h$: for large $h$ it equals the sum of the masses of two numerical semigroups attached to $A_0$, where the mass adds $L(L+1)/2$ over the maximal runs of $L$ gaps. We describe how the total reaches this limit: it can rise to it, fall to it, or climb past it and return. When the monomial curve of $A_0$ is arithmetically Cohen--Macaulay, two Ap\'ery sets give every undershoot exactly, and partial sums of a fixed spectrum bound the total from below, with equality at every level exactly when two threshold sequences are similarly ordered; the total then rises concavely to its limit, reached exactly at the sumsets regularity. We characterize the similarly ordered four-element sets. For $\{0,a,b,2b\}$ with $0<a<b$ coprime, the three-distance theorem gives the total in explicit quadratic pieces. A non-Cohen--Macaulay set exceeds its limit at every failing level $h$ with $hd$ at least the sum of the two conductors; every such four-element set does so, and a five-element example never does. For arbitrary sets, finite corner data give the exact undershoot over any window.}

\keywords{sumsets, numerical semigroups, Ap\'ery sets, Cohen--Macaulay monomial curves, three-distance theorem, gaps}

\pacs[MSC Classification]{Primary 11B13, 20M14; Secondary 13H10, 11J71, 11D07}
